\begin{frame}{Time Series Basics: Stationarity Conditions}
    \begin{center}
    \begin{tikzpicture}[
        node distance=1.5cm and 2cm,
        box/.style={draw=blue!50, thick, rounded corners, fill=blue!5, text width=5cm, align=center, inner sep=10pt},
        title/.style={font=\bfseries\large, text=blue!80!black, yshift=5pt}
    ]
        % SSS Node
        \node[box] (sss) {
            \node[title] at (0, 1.2) {Strictly Stationary (SSS)};
            \vspace{0.5cm}
            Joint distribution is invariant to time shifts $\tau$:
            $$F_X(x_{t_1}, \dots, x_{t_k}) = F_X(x_{t_1+\tau}, \dots, x_{t_k+\tau})$$
            \textit{(Too strict for real-world empirical data)}
        };

        % WSS Node
        \node[box, right=of sss, draw=green!60!black, fill=green!5] (wss) {
            \node[title, text=green!60!black] at (0, 1.2) {Weakly Stationary (WSS)};
            \vspace{0.5cm}
            First two moments are time-invariant:
            $$ \begin{cases}
            \mu_X(t) = \text{const}, \forall t \\
            R_X(t_i, t_j) = R_X(|t_i - t_j|)
            \end{cases} $$
            \textit{(Our primary assumption for modeling)}
        };

        % Arrow connecting them
        \draw[-{Stealth[scale=1.5]}, thick, gray, dashed] (sss) -- node[above, font=\footnotesize] {Implies} (wss);
    \end{tikzpicture}
    \end{center}
\end{frame}